\documentclass{beamer}
\usetheme{Madrid}
\usecolortheme{default}
\usepackage{amsmath, amssymb, graphicx, booktabs}
\usepackage{hyperref}

\title{GED Math Unit: Data, Statistics, and Probability}

\date{\today}

\begin{document}

% Title Slide
\begin{frame}
  \titlepage
\end{frame}

% Outline
\begin{frame}{Lesson Overview}
  \tableofcontents
\end{frame}

% ==========================
% LESSON 1: DATA REPRESENTATION
% ==========================

\section{Lesson 1: Data Representation}

\begin{frame}{Understanding Data}
\begin{itemize}
  \item Data: information collected from observations or measurements.
  \item Types:
  \begin{itemize}
    \item Qualitative (categorical)
    \item Quantitative (numerical)
  \end{itemize}
  \item Quantitative data can be:
  \begin{itemize}
    \item Discrete (countable values)
    \item Continuous (measured values)
  \end{itemize}
\end{itemize}
\end{frame}

\begin{frame}{Visual Representations}
  \begin{itemize}
    \item Bar Graphs -- Compare categories.
    \item Histograms -- Show frequency of numerical data.
    \item Line Graphs -- Show changes over time.
    \item Pie Charts -- Show proportions.
  \end{itemize}
  \begin{center}
%    \includegraphics[width=0.6\linewidth]{example-image-a}
  \end{center}
\end{frame}

\begin{frame}{Example 1: Bar Graph}
A survey found how many students prefer each subject:
\\
Math: 12, Science: 9, English: 15, History: 6.
\\
\textbf{Task:} Draw a bar graph.

\begin{center}
%\includegraphics[width=0.7\linewidth]{example-image-b}
\end{center}
\end{frame}

\begin{frame}{Example 2: Histogram}
Data: Test scores -- 45, 55, 60, 65, 70, 70, 75, 80, 85.
\\
\textbf{Task:} Create frequency intervals of 10.

\begin{tabular}{ll}
Interval & Frequency \\
40-49 & 1 \\
50-59 & 1 \\
60-69 & 2 \\
70-79 & 3 \\
80-89 & 2 \\
\end{tabular}
\end{frame}

\begin{frame}{Example 3: Line Graph}
A student’s GPA over 5 semesters: 2.8, 3.0, 3.2, 3.5, 3.6.
\\
\textbf{Task:} Plot points and connect them with a line.
\\
%\includegraphics[width=0.7\linewidth]{example-image-c}
\end{frame}

\begin{frame}{Example 4: Pie Chart}
Total students: 30
\\
Math: 12, Science: 9, English: 9.
\\
\textbf{Angles:}
\[
\text{Math: } \frac{12}{30} \times 360 = 144^\circ, \; \text{Science: } 108^\circ, \; \text{English: } 108^\circ
\]
%\includegraphics[width=0.6\linewidth]{example-image}
\end{frame}

\begin{frame}{Example 5: Stem-and-Leaf Plot}
Data: 42, 45, 46, 51, 53, 53, 56, 61.
\\
Stem | Leaf \\
4 | 2 5 6 \\
5 | 1 3 3 6 \\
6 | 1 \\
\end{frame}

% ==========================
% LESSON 2: MEASURES OF CENTRAL TENDENCY
% ==========================

\section{Lesson 2: Measures of Central Tendency}

\begin{frame}{Definitions}
\begin{itemize}
  \item Mean: Average value
  \item Median: Middle value when ordered
  \item Mode: Most frequent value
  \item Range: Difference between max and min
\end{itemize}
\end{frame}

\begin{frame}{Example 1: Mean}
Data: 5, 7, 9, 10, 15
\\
\[
\text{Mean} = \frac{5+7+9+10+15}{5} = 9.2
\]
\end{frame}

\begin{frame}{Example 2: Median}
Data: 3, 5, 7, 9, 11, 13
\\
\text{Median} = average of 7 and 9 = 8
\end{frame}

\begin{frame}{Example 3: Mode}
Data: 1, 2, 2, 3, 4, 4, 4, 5
\\
\text{Mode} = 4 (appears most often)
\end{frame}

\begin{frame}{Example 4: Range}
Data: 12, 15, 17, 19, 25
\\
\text{Range} = 25 - 12 = 13
\end{frame}

\begin{frame}{Example 5: Combined Measures}
Data: 8, 9, 9, 10, 11, 12
\\
Mean = 9.83, Median = 9.5, Mode = 9, Range = 4.
\end{frame}

% ==========================
% LESSON 3: PROBABILITY BASICS
% ==========================

\section{Lesson 3: Probability Basics}

\begin{frame}{Understanding Probability}
\begin{itemize}
  \item Probability measures how likely an event is.
  \item Formula:
  \[
  P(E) = \frac{\text{Number of favorable outcomes}}{\text{Total outcomes}}
  \]
  \item Range: 0 (impossible) to 1 (certain)
\end{itemize}
\end{frame}

\begin{frame}{Example 1: Coin Toss}
P(Heads) = \(\frac{1}{2}\)
\end{frame}

\begin{frame}{Example 2: Rolling a Die}
P(rolling 4) = \(\frac{1}{6}\)
\end{frame}

\begin{frame}{Example 3: Drawing Cards}
Deck of 52 cards.
\\
P(drawing a heart) = \(\frac{13}{52} = \frac{1}{4}\)
\end{frame}

\begin{frame}{Example 4: Complementary Events}
P(not rolling 6) = 1 - P(rolling 6) = 1 - \(\frac{1}{6} = \frac{5}{6}\)
\end{frame}

\begin{frame}{Example 5: Multiple Events}
P(A and B) = P(A) \(\times\) P(B) if independent.
\\
Example: Tossing 2 heads = \(\frac{1}{2}\times\frac{1}{2}=\frac{1}{4}\)
\end{frame}

% ==========================
% LESSON 4: ADVANCED PROBABILITY
% ==========================

\section{Lesson 4: Advanced Probability}

\begin{frame}{Dependent Events}
P(A and B) = P(A) \(\times\) P(B|A)\\
\vspace{0.2 cm}
Example: Drawing 2 red cards without replacement.
\\
\vspace{0.2 cm}
P(1st red) = \(\frac{26}{52}\), P(2nd red) = \(\frac{25}{51}\)\\
\vspace{0.2 cm}

So, P = \(\frac{26}{52}\times\frac{25}{51}=\frac{25}{102}\)
\end{frame}

\begin{frame}{Example 1: OR Events}
P(A or B) = P(A) + P(B) - P(A and B)
\\
\vspace{0.2 cm}
Example: Rolling a 2 or even number.
\\
\vspace{0.2 cm}
P(2)=\(\frac{1}{6}\), P(even)=\(\frac{3}{6}\), P(both)=\(\frac{1}{6}\)
\\
\vspace{0.2 cm}
P=\(\frac{3}{6}\)
\end{frame}

\begin{frame}{Example 2: Conditional Probability}
P(B|A)=\(\frac{P(A\text{ and }B)}{P(A)}\)
\\
\vspace{0.2 cm}
See these examples:\\
\vspace{0.2 cm}
\url{https://probability.oer.math.uconn.edu/wp-content/uploads/sites/2187/2018/01/prob3160ch4.pdf}
\end{frame}

\begin{frame}{Example 3: Tree Diagrams}

\begin{itemize}
    \item Helps us to visualize multiple outcomes.
    \item Usually drawn from left to right and shows ALL possible outcomes.
    \item Can be used for counting outcomes to find probability.
\end{itemize}

%\includegraphics[width=0.6\linewidth]{example-image}
\end{frame}

\begin{frame}{Example 4: Expected Value}
If probabilities and outcomes are known:
\vspace{0.2 cm}
\\
\(E = \sum x_i p_i\)
\\
\vspace{0.2 cm}
Example: The expected value is the average outcome of a random event over many trials, calculated by multiplying the value of each outcome by its probability and summing the results. For example, a raffle ticket that costs \$3 and has a 1/150 chance of winning a \$300 prize has an expected value of -\$1. \\
This is because the expected value is \\
\vspace{1.0 cm}
$\left(297 * \frac{1}{150}\right) + \left(-3 * \frac{149}{150}\right),$ which equals approximately -\$1. 
\end{frame}

\begin{frame}{Example 5: Permutations/Combinations}
Permutations: order matters. \(nP_r=\frac{n!}{(n-r)!}\)
\\
\vspace{0.2 cm}
Combinations: order doesn’t matter. \(nC_r=\frac{n!}{r!(n-r)!}\)
\\
\vspace{0.2 cm}
Combination Example: Choose 3 students from 5: \(5C3=10\)\\
\vspace{0.2 cm}

A simple permutation example is arranging the letters 'a', 'b', and 'c', where the different orders like abc, acb, bac, bca, cab, and cba are all distinct permutations. 

The letters a, b, and c can be arranged in 6 different ways: 
abc
acb
bac
bca
cab
cba

\end{frame}

% ==========================
% LESSON 5: INTERPRETING DATA
% ==========================

\section{Lesson 5: Interpreting Data}

\begin{frame}{Trends and Outliers}
\begin{itemize}
  \item Outlier: data point far from others.
  \item Trends: upward/downward patterns.
  \item Use scatter plots and best-fit lines.
\end{itemize}
%\includegraphics[width=0.7\linewidth]{example-image-c}
\end{frame}

\begin{frame}{Example 1: Scatter Plot}
Height vs. Weight: correlation = positive.
\\
%\includegraphics[width=0.6\linewidth]{example-image}
\end{frame}

\begin{frame}{Example 2: Line of Best Fit}
Predict value from trend.
\\
Equation: \(y = mx + b\)
\end{frame}

\begin{frame}{Example 3: Misleading Graphs}
Check for:
\begin{itemize}
  \item Truncated axes
  \item Distorted scales
  \item Selective data
\end{itemize}
\end{frame}

\begin{frame}{Example 4: Two-Way Tables}
\begin{tabular}{lcc|c}
 & Male & Female & Total \\
\hline
Passed & 15 & 20 & 35 \\
Failed & 5 & 10 & 15 \\
\hline
Total & 20 & 30 & 50 \\
\end{tabular}
\\
P(Female|Passed)=\(\frac{20}{35}=0.57\)
\end{frame}

\begin{frame}{Example 5: Interpreting Mean and Spread}
\textbf{Standard Deviation:} measures spread.
\\
Low SD: data clustered, High SD: data spread out.
\\
Example: Test A SD=2, Test B SD=10 → Test B more variable.
\end{frame}

\begin{frame}{Summary}
\begin{itemize}
  \item Representing data visually
  \item Central tendency and spread
  \item Probability and combinations
  \item Interpreting statistical results
\end{itemize}
\end{frame}

\end{document}
